\subsection{Arithmetic Kummer coordinates in the dyadic action kernel}

% Arithmetic pro-two-kernel Kummer coordinates.
% Required bibliography key: HainMatsumoto2003.
\begin{proposition}[Arithmetic Kummer coordinates in the pro-two kernel]
\label{prop:arithmetic-kummer-kernel-axes}
Let \(G=\widehat{\mathrm{GT}}\), let \(H\) be its arithmetic subgroup, and put
\[
 M=\ker\lambda,\qquad N=\ker(\lambda,\rho_2^I).
\]
\[
 D=\bigcap_{q\ {\rm prime}}\ker(\rho_q^I|_M),
\]
\[
 R_M=\{g\in M:f_g=1\text{ in the free pro-two group}\}.
\]
For every odd rational prime \(q\), there is an element
\(k_q\in R_M\cap N\cap H\) whose only nonzero rational Kummer
coordinate is
\[
 \rho_{q,2}^I(k_q)=-2.
\]
In particular the tuple of dyadic odd-prime coordinates on the closed
arithmetic subgroup of \(R_M\cap N\) killing every odd-primary
Kummer coordinate has image exactly
\begin{equation}
 \prod_{q\ {\rm odd}}2\mathbf Z_2.                       \label{eq:arithmetic-kernel-axes-1}
\end{equation}
If \(G=DH\), and \(R_D=R_M\cap D\), the natural map
\[
 (R_M\cap N)/R_D\longrightarrow N/D
\]
is an isomorphism. If, in addition,
\(M=R_M(M\cap H)\), then
\[
 D=R_D(D\cap H),\qquad G=R_DH.
\]
\end{proposition}

\begin{proof}
Let \(\varepsilon_2\in\widehat{\mathbf Z}\) have dyadic component
one and every odd-primary component zero.
The tangential action on the pro-two fundamental group of
\(\mathbf P^1-\{0,1,\infty\}\) is unramified at every odd prime:
apply \cite[\S8.5, p.~151]{HainMatsumoto2003} with \(N=1\) and
\(\ell=2\). This is the action into automorphisms, so no inner
automorphism ambiguity is involved.

Choose an inertia element \(\tau_q\) with dyadic Kummer coordinate
one at \(q\). It has coordinate zero at every other rational
prime: these primes are \(q\)-adic units, whose two-power roots
lie in unramified extensions, as do the two-power roots of unity.
Choose complex conjugation \(c\) fixing the compatible positive
real roots of the positive rational primes. Its Kummer coordinates
are zero and its multiplier is \(-1\). Then
\[
 k_q=[c,\tau_q]^{\varepsilon_2}
\]
has full multiplier one and trivial pro-two action.
The crossed Kummer law gives its dyadic coordinate \(-2\) at \(q\)
and zero elsewhere, while the power \(\varepsilon_2\) kills all
odd-primary coordinates. This proves the asserted individual
coordinates.

An arithmetic element of \(M\) fixes every rational square root,
because the quadratic fields are cyclotomic. Its dyadic coordinates
are therefore even. The image in \eqref{eq:arithmetic-kernel-axes-1} is closed by compactness,
contains each coordinate axis and every \(\mathbf Z_2\)-multiple
of that axis, and hence contains the closure of their direct sum.
This proves \eqref{eq:arithmetic-kernel-axes-1}, including its simultaneous assertion.

The normalized pro-two action has pro-two image: its actions on
the finite algebras
\(\mathbf F_2\langle S,T\rangle/(S,T)^d\) are upper unitriangular
and detect its parameter. Consequently
\begin{equation}
 a^{\,1-\varepsilon_2}\in R_M,\qquad
 a^{\varepsilon_2}R_M=aR_M\qquad(a\in M).                \label{eq:arithmetic-kernel-axes-2}
\end{equation}
These powers are taken in the procyclic closure of \(a\).

Suppose \(G=DH\). For \(n\in N\), choose \(a\in N\cap H\)
with \(na^{-1}\in D\). The element \(a^{1-\varepsilon_2}\)
lies in \(R_M\cap N\cap H\) and realizes all its odd-primary
coordinates. Formula \eqref{eq:arithmetic-kernel-axes-1} realizes the remaining dyadic
odd-prime tuple inside the same group, without changing any
other coordinate. Thus \(R_M\cap N\to N/D\) is surjective.
Its kernel is \(R_M\cap D=R_D\), and the compact quotient theorem
gives the stated isomorphism.

Finally take \(d\in D\), and suppose
\(M=R_M(M\cap H)\). Choose \(a\in M\cap H\) with the same
pro-two parameter as \(d\). The dyadic beta measure, and hence
its mass \(\rho_{2,2}^I/2\), depends only on that parameter;
thus \(\rho_{2,2}^I(a)=0\). Replace \(a\) by
\(a^{\varepsilon_2}\), preserving its parameter by \eqref{eq:arithmetic-kernel-axes-2}
and killing all its odd-primary Kummer coordinates.
Use \eqref{eq:arithmetic-kernel-axes-1} to cancel its even dyadic coordinates at every odd
prime by an element \(k\in R_M\cap N\cap H\).
Then \(a^{\varepsilon_2}k^{-1}\in D\cap H\) still matches \(d\).
This gives \(D=R_D(D\cap H)\), and \(G=DH\) gives \(G=R_DH\).
\end{proof}
